\documentclass[12pt,a4paper]{article}

% ── Packages ──────────────────────────────────────────────────────────────────
\usepackage[margin=2.5cm]{geometry}
\usepackage{amsmath, amssymb}
\usepackage{graphicx}
\usepackage{booktabs}
\usepackage{hyperref}
\usepackage{cleveref}
\usepackage{caption}
\usepackage{subcaption}
\usepackage{listings}
\usepackage{xcolor}
\usepackage{setspace}
\usepackage{parskip}
\usepackage{fancyhdr}
\usepackage{titlesec}
\usepackage{biblatex}
\addbibresource{references.bib}

% ── Page Style ────────────────────────────────────────────────────────────────
\pagestyle{fancy}
\fancyhf{}
\rhead{4006 Task — Flight Dynamics Simulation}
\lhead{\leftmark}
\cfoot{\thepage}

% ── Code Listing Style ────────────────────────────────────────────────────────
\lstdefinestyle{matlab}{
    language=Matlab,
    basicstyle=\ttfamily\small,
    keywordstyle=\color{blue},
    commentstyle=\color{gray},
    stringstyle=\color{orange},
    numbers=left,
    numberstyle=\tiny\color{gray},
    breaklines=true,
    frame=single,
    captionpos=b,
}

% ── Title ─────────────────────────────────────────────────────────────────────
\title{
    \vspace{2cm}
    \textbf{Flight Dynamics Simulation Model} \\[0.5em]
    \large 4006 Task Report \\[0.5em]
    \normalsize Department of Aerospace Engineering
}
\author{Your Name \\ Student ID: XXXXXXX \and Second Author \\ Student ID: XXXXXXX}
\date{\today}

% ══════════════════════════════════════════════════════════════════════════════
\begin{document}
% ══════════════════════════════════════════════════════════════════════════════

\maketitle
\thispagestyle{empty}
\newpage

\tableofcontents
\newpage

% ── 1. Introduction ───────────────────────────────────────────────────────────
\section{Introduction}
\label{sec:intro}

Provide background on the project. State the objectives and outline the
structure of the report.

% ── 2. Forces and Torques Model ───────────────────────────────────────────────
\section{Forces and Torques Model}
\label{sec:forces}

\subsection{Aerodynamic Forces}

Describe lift, drag, and side force equations.

\begin{equation}
    \mathbf{F}_{aero} = \begin{bmatrix} F_x \\ F_y \\ F_z \end{bmatrix}
    \label{eq:aero_forces}
\end{equation}

\subsection{Aerodynamic Moments}

Describe rolling, pitching, and yawing moments.

\begin{equation}
    \mathbf{M}_{aero} = \begin{bmatrix} l \\ m \\ n \end{bmatrix}
    \label{eq:aero_moments}
\end{equation}

\subsection{Gravitational Forces}

Describe gravitational force components in the body frame.

% ── 3. Control Surface Model ──────────────────────────────────────────────────
\section{Control Surface Model}
\label{sec:control}

\subsection{Control Inputs}

Describe aileron ($\delta_a$), elevator ($\delta_e$), and rudder ($\delta_r$)
deflections.

\subsection{Aerodynamic Derivatives}

Explain control surface aerodynamic derivatives and their effect on forces and
moments.

% ── 4. Trim Analysis ──────────────────────────────────────────────────────────
\section{Trim Analysis}
\label{sec:trim}

\subsection{Trim Conditions}

Define the equilibrium (trim) conditions for steady-level flight.

\subsection{Trim Solution}

Describe the numerical method used to solve for trim states and control
deflections.

\subsection{Trim Validation}

Present validation results.

\begin{table}[h!]
    \centering
    \caption{Trim State Results}
    \label{tab:trim}
    \begin{tabular}{lcc}
        \toprule
        \textbf{Variable} & \textbf{Symbol} & \textbf{Value} \\
        \midrule
        Airspeed          & $V_a$           & -- m/s         \\
        Angle of attack   & $\alpha$        & -- deg         \\
        Pitch angle       & $\theta$        & -- deg         \\
        Elevator          & $\delta_e$      & -- deg         \\
        Throttle          & $\delta_t$      & --             \\
        \bottomrule
    \end{tabular}
\end{table}

% ── 5. Autopilot Design ───────────────────────────────────────────────────────
\section{Autopilot Design}
\label{sec:autopilot}

\subsection{Autopilot Architecture}

Describe the overall autopilot structure (inner/outer loops, etc.).

\subsection{Gain Computation}

Explain how autopilot gains were computed (refer to
\texttt{compute\_autopilot\_gains.m}).

\subsection{Simulation Results}

Present and discuss simulation results. Include figures as needed.

\begin{figure}[h!]
    \centering
    % \includegraphics[width=0.8\textwidth]{figures/your_figure.png}
    \caption{Autopilot simulation response.}
    \label{fig:autopilot_response}
\end{figure}

% ── 6. Discussion ─────────────────────────────────────────────────────────────
\section{Discussion}
\label{sec:discussion}

Interpret your results. Compare against expectations or reference values.
Discuss any discrepancies or limitations.

% ── 7. Conclusion ─────────────────────────────────────────────────────────────
\section{Conclusion}
\label{sec:conclusion}

Summarise findings and their significance. Suggest future work if appropriate.

% ── References ────────────────────────────────────────────────────────────────
\newpage
\printbibliography

% ── Appendix ──────────────────────────────────────────────────────────────────
\newpage
\appendix

\section{MATLAB Code}
\label{app:code}

\subsection{forces\_moments.m}
\begin{lstlisting}[style=matlab, caption={Forces and moments computation}]
% Paste or \input your MATLAB code here
\end{lstlisting}

\subsection{compute\_trim.m}
\begin{lstlisting}[style=matlab, caption={Trim computation}]
% Paste or \input your MATLAB code here
\end{lstlisting}

\end{document}
